\begin{theorem}[Valuation of the wild odd correction]\label{mod:cases-wildodd-correctionvaluation}
Let $d$ be wild odd correction data, and put
\[
 p=\operatorname{char}k_F=[K:F],\qquad
 a=T-m,\qquad n=N_{K/F}(u),
\]
\[
 e_0=\left\lfloor\frac{(p-1)T}{p}\right\rfloor
     =T-\left\lceil\frac Tp\right\rceil,
 \qquad
 B=s_1+\sum_{q\in[1,p-1)}(-n)^q s_{p-q}.
\]
Here all quotients in the definitions of $e_0$ and the auxiliary depths
are integer floors.  The identity with the integer ceiling is the exact
one supplied by the symmetric-bound API, and for odd $p$ and $T\ge2$ that
API gives $T\le2e_0$.

Lean first proves the residue-prime valuation statements in the two
characteristic branches.  If the field element $(p:F)$ is zero, then
\[
 v_F(p)=\top
\]
exactly.  If $(p:F)\ne0$, there is an integer $e$ with
\[
 v_F(p)=e,\qquad e_0\le e.
\]
Thus the mixed-characteristic ramification input is the lower bound
$e\ge e_0$, not the generally false equality $e=e_0$, and
$T\le e+e_0$.

Write $b=m-1$ and $t=T-1$, as in the manuscript.  Since $a=t-b$, the Lean
proof uses the following exhaustive, pairwise disjoint trichotomy without
weakening either strict boundary:
\[
 b<t\ \Longleftrightarrow\ m<T\ \Longleftrightarrow\ 0<a,
\quad
 b=t\ \Longleftrightarrow\ m=T\ \Longleftrightarrow\ a=0,
\quad
 t<b\ \Longleftrightarrow\ T<m\ \Longleftrightarrow\ a<0.
\]
Before estimating the quotient in the closed correction formula, Lean
proves the exact denominator valuation separately in all three cases:
\[
 v_F(1-n^{p-1})=
 \begin{cases}
  0,&0<a,\\
  0,&a=0,\\
  (p-1)a,&a<0.
 \end{cases}
\]
The boundary equality uses the supplied noncancellation of every
$n+[j]$ and the exact Teichm\"uller product; it is not inferred by merging
$a=0$ with a weak positive case.  Each displayed valuation is finite, so
the denominator is nonzero before quotient valuation and cancellation are
used.

For $a\ge0$, the symmetric-function estimate is applied to $s_1$ and,
for every $q\in[1,p-1)$, to the exact index $p-q$, whose range is
$2\le p-q<p$.  It gives
\[
 v_F(s_j(u))\ge
 \left\lfloor\frac{ja+(p-1)T}{p}\right\rfloor\ge e_0
 \quad(1\le j<p),
 \qquad v_F(B)\ge e_0.
\]
The powers $(-n)^q$ have nonnegative valuation.  Together with the
separately proved unit denominator and $v_F(p)=e$, this yields
$v_F(X)\ge e+e_0\ge T$ in both the strict $a>0$ and boundary $a=0$
branches.

For $a<0$, put $c=-a>0$ and $y=u^{-1}$.  Lean proves the finite-family
reciprocal identity
\[
 s_j(u)=n\,s_{p-j}(y)\qquad(1\le j\le p-1)
\]
and applies the symmetric bound to $y$ with
\[
 L_r(c)=\left\lfloor\frac{rc+(p-1)T}{p}\right\rfloor.
\]
After subtracting the exact denominator valuation $(p-1)a$, the trace term
uses the endpoint $r=p-1$ and the remaining finite sum uses precisely
$2\le j<p$:
\[
 e_0\le(p-2)c+L_{p-1}(c),
\qquad
 e_0\le(j-2)c+L_{p-j}(c)\quad(2\le j<p).
\]
These are the literal reciprocal floor inequalities from SymmetricBounds;
no endpoint $j=p$ is introduced.  Termwise valuation bounds followed by
the finite-sum inequality give
\[
 v_F(B)\ge e_0+(p-1)a,
\]
and cancellation against the denominator proves the same unweakened final
depth in the negative branch.  In equal characteristic the exact closed
identity instead has $(p:F)=0$ and proves $X=0$.  In all cases the primary
Lean theorem is therefore
\[
 \boxed{\quad X\in\mathfrak p_F^T.\quad}
\]

Finally, let $\tau$ be packaged exact-conductor data with
$m_F(\tau)=T$, let $0<r\le T$, and suppose the stationary linearization
\[
 \tau(1+x)=\psi(Ax)\qquad(x\in\mathfrak p_F^r)
\]
is available.  Lean derives, rather than assumes, the full additive
triviality statement
\[
 x\in\mathfrak p_F^T\quad\Longrightarrow\quad\psi(Ax)=1
\]
from exact conductor triviality on $U_F^T$ and the inclusion
$\mathfrak p_F^T\subseteq\mathfrak p_F^r$.  Applying this result to
$-X\in\mathfrak p_F^T$ gives the exact phase
\[
 \boxed{\quad\psi(-AX)=1.\quad}
\]
This is the lattice and phase interface exposed for the later wild-odd
main theorem; no parity reduction, phase reduction, or First Main Lemma is
used.
\LeanFile{LanglandsFirstMainLemma/Cases/WildOdd/CorrectionValuation.lean}
\lean{LanglandsFirstMainLemma.wildOdd_correction_mem}
\lean{LanglandsFirstMainLemma.WildOddCorrectionData.ord_one_sub_n_pow_of_pos}
\lean{LanglandsFirstMainLemma.WildOddCorrectionData.ord_one_sub_n_pow_of_eq_zero}
\lean{LanglandsFirstMainLemma.WildOddCorrectionData.ord_one_sub_n_pow_of_neg}
\lean{LanglandsFirstMainLemma.wildOdd_prime_ord_lower_bound}
\lean{LanglandsFirstMainLemma.wildOdd_prime_ord_eq_top_of_eq_zero}
\lean{LanglandsFirstMainLemma.wildOdd_prime_ord_exists_of_ne_zero}
\lean{LanglandsFirstMainLemma.wildOdd_additive_triviality}
\lean{LanglandsFirstMainLemma.wildOdd_correction_phase}
\leanok
\SourceText{Lemma lem:p-valuation and Proposition prop:X-depth.}
\uses{mod:cases-wildodd-correctionformula,mod:ramification-symmetricbounds,mod:ramification-pullbackconductors}
\FormalizationContract The closed formula is used as an exact field
identity.  The three break-position cases remain disjoint; strict
inequalities, the boundary noncancellation argument, every floor and
ceiling conversion, denominator valuation, reciprocal index range, and
the final lattice exponent $T$ are retained literally.  The phase theorem
derives additive triviality from exact multiplicative conductor data and
stationary linearization before applying it to the proved correction
depth.  No helper assumes the desired depth, and no downstream wild-odd or
First Main Lemma theorem is imported.
\end{theorem}
